\begingroup%
\definecolor{ksBlue}{HTML}{3478B8}%
\definecolor{ksTeal}{HTML}{27877F}%
\definecolor{ksAmber}{HTML}{C58B28}%
\definecolor{ksInk}{HTML}{243447}%
\definecolor{ksPaleBlue}{HTML}{F0F6FC}%
\definecolor{ksPaleTeal}{HTML}{EFF8F5}%
\definecolor{ksPaleAmber}{HTML}{FFF8E8}%
\begin{tikzpicture}[
  x=1mm, y=1mm,
  card/.style={draw, fill=white, rounded corners=2mm, line width=.7pt,
    text=ksInk, align=flush center, font=\fontsize{9.5}{12.5}\selectfont,
    text width=69mm, minimum width=73mm, minimum height=14mm, inner xsep=2mm, inner ysep=1.5mm,
    outer sep=.35pt},
  challenge/.style={card, font=\bfseries\fontsize{9.5}{12.5}\selectfont},
  frame/.style={draw, rounded corners=3mm, line width=.7pt, inner sep=3mm,
    outer sep=.35pt},
  heading/.style={text=ksInk, font=\bfseries\fontsize{10}{12.5}\selectfont,
    inner sep=0pt, outer sep=0pt},
  link/.style={line width=.65pt},
  flow/.style={-{Stealth[length=2.6mm,width=1.9mm]}, ksInk, line width=.9pt},
  edge text/.style={text=ksInk, font=\fontsize{9}{12.5}\selectfont,
    inner sep=1mm}
]
  \node[heading] (s1-title) at (40,0) {环节一：数据与表示};
  \node[challenge, draw=ksBlue] (c1-title) at (0,-12)
    {挑战一：硬件感知不足};
  \node[card, draw=ksBlue] (c1-a) at (80,-12)
    {语料偏置：以 NVIDIA 旧架构为主\\
    Blackwell 与国产卡语料稀缺};
  \node[card, draw=ksBlue] (c1-b) at (80,-29)
    {语言杂（语言是思维的边界）：\\
    CUDA / Triton / TileLang / Ascend 等};
  \node[card, draw=ksBlue] (c1-c) at (0,-29)
    {内容杂（源码未经压缩与蒸馏）：\\
    Host / IO 样板淹没设计决策};

  \node[heading] (s2-title) at (40,-51) {环节二：生成与迭代};
  \node[challenge, draw=ksTeal] (c2-title) at (0,-63)
    {挑战二：迭代收敛困难};
  \node[card, draw=ksTeal] (c2-a) at (80,-63)
    {随机性高、一致性低：\\
    同条件重复运行，轨迹明显分化};
  \node[card, draw=ksTeal] (c2-b) at (80,-80)
    {设计与实现强耦合：\\
    好设计因实现瑕疵被“误杀”};
  \node[card, draw=ksTeal] (c2-c) at (0,-80)
    {避难就易走弯路：为尽快跑通\\
    退回 mma.sync 等旧指令，平台期长};

  \node[heading] (s3-title) at (40,-102) {环节三：交付与维护};
  \node[challenge, draw=ksAmber] (c3-title) at (0,-114)
    {挑战三：维护实用麻烦};
  \node[card, draw=ksAmber] (c3-a) at (80,-114)
    {Shape 特化分支膨胀：\\
    以 if-else 针对特定负载“刷分”，缺乏泛化};
  \node[card, draw=ksAmber] (c3-b) at (80,-131)
    {隐式缺陷：共享内存竞争、越界、死锁\\
    测试集随之膨胀，迭代变慢};
  \node[card, draw=ksAmber] (c3-c) at (0,-131)
    {黑盒代码难以审阅：\\
    设计取舍埋在底层代码中，上线风险高};

  \begin{scope}[on background layer]
    \node[frame, draw=ksBlue, fill=ksPaleBlue,
      fit=(s1-title)(c1-title)(c1-a)(c1-b)(c1-c)] (s1) {};
    \node[frame, draw=ksTeal, fill=ksPaleTeal,
      fit=(s2-title)(c2-title)(c2-a)(c2-b)(c2-c)] (s2) {};
    \node[frame, draw=ksAmber, fill=ksPaleAmber,
      fit=(s3-title)(c3-title)(c3-a)(c3-b)(c3-c)] (s3) {};
  \end{scope}

  \draw[link, ksBlue] (c1-title) -- (c1-a);
  \draw[link, ksBlue] (c1-a) -- (c1-b);
  \draw[link, ksBlue] (c1-b) -- (c1-c);
  \draw[link, ksTeal] (c2-title) -- (c2-a);
  \draw[link, ksTeal] (c2-a) -- (c2-b);
  \draw[link, ksTeal] (c2-b) -- (c2-c);
  \draw[link, ksAmber] (c3-title) -- (c3-a);
  \draw[link, ksAmber] (c3-a) -- (c3-b);
  \draw[link, ksAmber] (c3-b) -- (c3-c);
  \draw[flow] (s1.south) --
    node[edge text, right=1mm] {带着不完整的先验进入搜索} (s2.north);
  \draw[flow] (s2.south) --
    node[edge text, right=1mm] {产出高风险的黑盒代码} (s3.north);
\end{tikzpicture}%
\endgroup%
